\section{Expanded reasoning: alignment to a defensible writer profile}
\label{sec:expand-style}
\subsection{Body-band estimation: what the horizontal profile guesses}
Let $M$ be a binary line mask and $p(y)=\sum_x M(y,x)$ its horizontal ink
projection. The body-band heuristic takes the first and last rows above
$0.45\max_y p(y)$ as top and baseline $b$, with $h_x=b-y_{\rm top}$.
The baseline here is an image-derived estimate, not an anatomical fact about
every glyph. Capitals, long descenders, isolated punctuation and uneven ink
can bias the threshold. Bands narrower than three pixels are rejected because
normalised measurements would become unstable. A 45\% profile threshold is
different from a threshold on an individual pixel's brightness.

Constructed profile $(0,2,8,10,9,3,0)$ has cutoff 4.5, so rows 2 through 4
define the dense band. Using half-open row edges would give a different
height by one pixel; the convention must follow the implementation. Changing
0.45 changes estimated x-height, and therefore changes all later quantities
expressed in x-heights. A percentile of ink rows, fitted baseline, or
character-specific estimates are alternatives, with different assumptions.
The exact 0.45 and three-pixel guard are project heuristics~\cite{code2026}.

\subsection{Refining shared character cuts: every term and its purpose}
CTC alignment gives approximate character spans. Their adjacent centres
provide initial shared cuts. For column profile $p_x$ smoothed over a short
window, the report uses
\[
C_x={p_x\over q_{0.60}(p_j:p_j>0)}+
1.15\operatorname{clip}\left({b-y_{\rm top}(x)\over h_x},0,3\right).
\]
The first term measures how much ink a cut would cross relative to a typical
nonempty column. Excluding zero columns avoids a zero scale on sparse text;
the numerical denominator floor handles degenerate cases. The second term
penalises columns containing tall ink, favouring low connectors near the
baseline over stems. Clip bounds limit the influence of extreme ascenders.
The two terms are dimensionless; 1.15 sets their relative priority.

For smoothed ink count 2, positive-column percentile 8, height above baseline
4 and $h_x=20$, the terms are 0.25 and $1.15(0.2)=0.23$, total 0.48.
A taller column with the same count and height 20 scores 1.40. A true blank
column is explicitly assigned zero cost. This is a constructed arithmetic
example, not a measured cut. There is no published reason why 1.15 rather
than 1.0 is optimal for these writers.

Read-only inspection also shows a smoothing width near $0.12h_x$, snapped to
an odd integer of at least three; a search radius near $0.5h_x$ of at least
three; and a proximity penalty $0.6|x-x_{\rm CTC}|/r$ inside that search.
This last 0.6 is dimensionless and is unrelated to the 0.6 source-pixel
polyline tolerance or the 0.6 mm hardware join. Ordered boundaries must
leave $\max(1,0.08h_x)$ pixels per cell. At $h_x=25$ that minimum is 2
pixels. Without ordering, two neighbouring boundaries can snap to the same
gap and create an empty character. Increasing minimum width may instead
force a cut into a narrow genuine glyph. These are interacting heuristics;
a local cost minimum alone cannot establish correct character ownership.

\subsection{Binarisation, square dilation and Zhang--Suen thinning}
Binarisation maps grey values to an ink/not-ink decision. It discards
intensity information; later thinning cannot restore a faint stroke lost at
this step. Dilation with square radius $r$ sets a pixel if any mask pixel
within $\max(|\Delta x|,|\Delta y|)\le r$ is on. For radius one, an isolated
pixel becomes a $3\times3$ block. Erosion requires all neighbourhood pixels
to be on; closing combines the two in order. Boundary padding changes their
behaviour at the edge.

For Zhang--Suen thinning, label the clockwise neighbours
$p_2=N,p_3=NE,p_4=E,p_5=SE,p_6=S,p_7=SW,p_8=W,p_9=NW$.
Let $B=\sum_{i=2}^9p_i$, and $A$ be the number of 0-to-1 transitions in
that circular list. A foreground pixel is eligible when $2\le B\le6$
and $A=1$. In the first subiteration also require
$p_2p_4p_6=0$ and $p_4p_6p_8=0$; in the second require
$p_2p_4p_8=0$ and $p_2p_6p_8=0$. Delete marked pixels simultaneously
within each subiteration and repeat until stable~\cite{zhangsuen1984,zhanglink}.
The rules protect endpoints and local connectivity while eroding boundaries.
Sequential in-place deletion would implement a different algorithm.

The skeleton approximates stroke geometry; it does not reconstruct pen
pressure, velocity, original stroke order or an exactly equidistant medial
axis. Thinning the complete line before splitting cells avoids artificial
end conditions at every character boundary. Nearby strokes may still form
ambiguous junctions. A distance-transform medial axis and vector contour
fitting are alternatives, but they have different branch/noise behaviour.

\subsection{Graph tracing, junction continuation and detached marks}
Treat skeleton pixels as vertices and touching neighbours as edges. Degree
one suggests an endpoint; degree two suggests continuation; larger degree
indicates a junction. Short chains can be raster artefacts or real dots.
Therefore a length filter must be combined with location/character context.
At a junction the tracer chooses compatible continuations rather than
claiming the original hand followed that order.

For outward endpoint directions $u$ and $v$, straight continuation means
they point oppositely. A continuation deflection can be written
$\alpha=\arccos[(-u)\cdot v/(\|u\|\|v\|)]$, zero for straight passage.
The report's limit is 75 degrees. If directions were instead both measured
along travel, the minus sign would be omitted; mixing these conventions
reverses the test. The threshold prevents a sharp turn from being accepted
as a smooth continuation. It is permissive compared with 30 degrees and
restrictive compared with 90, but no universal optimum follows from geometry.
Detached dots are recovered as short strokes; crossbars require association
rules so they are not mistaken for inter-character connectors. There is
inherent ambiguity without online pen traces.

\subsection{Douglas--Peucker: bounded simplification, units and example}
Given endpoints $A,B$ and intermediate point $P$, point-to-segment distance is
\[
t=\operatorname{clip}\left({(P-A)\cdot(B-A)\over\|B-A\|^2},0,1\right),
\qquad d(P,AB)=\|P-[A+t(B-A)]\|.
\]
If $A=B$, use $\|P-A\|$. Find the farthest intermediate point. If its
distance exceeds $\epsilon$, retain it and recurse on the two subpaths;
otherwise replace that subpath by its chord. This is the
Douglas--Peucker principle~\cite{douglaspeucker1973,douglaslink}.
The inspected project uses the infinite supporting-line distance
$|\Delta x(P_y-A_y)-\Delta y(P_x-A_x)|/(\|B-A\|+10^{-9})$, rather than
the clipped-segment formula above. The distinction matters near endpoints:
distance to a supporting line can be small beyond a chord's endpoints.
Consequently the project's tolerance is not an unrestricted symmetric
Hausdorff bound between original and simplified curves.

For $(0,0),(1,0.2),(2,0)$, the intermediate distance is 0.2. At
$\epsilon=0.6$ it can be removed; at 0.1 it must be retained. At source
x-height 30 pixels, 0.6 pixels is $0.02$ x-heights; drawing at 4 mm per
x-height gives a nominal source-to-output scale of 0.08 mm. This is a
scale conversion of the simplification tolerance, not a bound on physical
pen error. Resampling, skeleton error, smoothing, joining and machine error
all add different distortions.

The report records 3738 vertices reduced to 974 at 0.6 source pixels,
leaving about 26.1\% and removing about 73.9\%. That demonstrates compression
for one line, not why 0.6 is optimal. Larger tolerances reduce G-code segments
and can erase small loops/corners; smaller ones preserve raster jitter and
increase work. Closed loops require care because identical endpoints make
the initial chord degenerate; the inspected routine treats near-closed paths
specially. Polygon topology and curvature are not generally guaranteed by
ordinary Douglas--Peucker simplification.

\subsection{Coordinate normalisation and shear: a worked transformation}
For source pixel $(x,y)$, left cut $x_L$, baseline $b$, estimated x-height
$h_x>0$, slant $\phi$ and trim $\delta$,
\[
v={b-y\over h_x},\qquad u={x-x_L\over h_x}-v\tan\phi-\delta.
\]
The first equation flips image-down to glyph-up. Division removes source
scale. Subtracting $v\tan\phi$ removes slant; subtracting $\delta$ trims
empty lead space, recorded separately so spacing can be restored. Shear is
an affine transformation, not rotation: it shifts horizontal position in
proportion to vertical height~\cite{wolberg1990}.

Take $x=40,y=30,x_L=20,b=50,h_x=20,\phi=10$ degrees and $\delta=0.1$.
Then $v=1$ and $u=1-0.17633-0.1\simeq0.72367$. Restoring lead 0.1
and the same shear at $H=4$ mm gives horizontal offset
$4(0.1+0.72367+0.17633)=4$ mm and vertical offset 4 mm.
If slant were applied both during extraction and placement without undoing
it, the letter would be sheared twice. A negative or tiny estimated x-height
invalidates all these units, hence explicit guards.

\subsection{Size gates: 2.1 and 2.9 are widths, not heights}
In the inspected source, ordinary glyph width is limited to 2.1 x-heights
and designated wide glyphs to 2.9~\cite{code2026}. The supplied list calls
these height/width factors; the report and named constants identify both
as \emph{width} limits. At $h_x=30$ pixels, they correspond to 63 and 87
pixels of normalised source span. They target bad cuts containing neighbouring
letters while allowing wider forms such as $m$ and $w$. They do not derive
from typography as universal ratios.

Other gates inspect height class, longest stroke, total length, stroke count,
advance and within-character distributions. Ascenders and descenders should
not be filtered by the same height expectation as a dot. Relative distributions
allow writer-specific scale, but small samples can make them unreliable.
Relaxing gates retains more personal variation and more fragments; tightening
them reduces corruption and can eliminate rare valid forms. Missing characters
then need explicit fallbacks whose provenance must remain visible.

\subsection{The $12\times12$ descriptor and two different distances}
The inspected shape-grid routine samples points along each segment, maps
the glyph's bounding box to a $12\times12$ grid, accumulates counts, takes
their square roots, applies a small blur, flattens to 144 values and normalises
to unit Euclidean length. Sampling avoids representing a long straight segment
by only its endpoints. Square roots reduce dominance by heavily sampled cells.
The blur weights are centre 0.25 and eight neighbours 0.09375 each, summing
to one. Flattening changes storage shape, not geometry.

Independent horizontal/vertical bounding-box scaling reduces translation and
size sensitivity but also removes aspect-ratio information. Width/height gates
therefore provide information the descriptor deliberately discards. Grid size
12 is a resolution/cost compromise: coarse grids alias details, fine grids
become more sensitive to alignment and sampling. It is a project choice.

The pooled class reference is a componentwise median of descriptors followed
by unit normalisation. For unit vectors $v_g,r_c$, prior distance is
$d=1-v_g^Tr_c$. Identical descriptors give zero. Because these descriptors
are nonnegative, their dot products lie in $[0,1]$, so this distance lies in
$[0,1]$ (up to rounding), rather than the $[0,2]$ range of arbitrary signed
unit vectors. A cosine distance need not obey the triangle inequality, so
``distance'' does not automatically mean a mathematical metric.

Within-writer consensus instead uses $D_{ij}=\|v_i-v_j\|_1$, the sum of
absolute coordinate differences. The medoid is
$i^*=\arg\min_i\sum_jD_{ij}$. With scalar examples $0,0.1,0.9$, the sums
are $1.0,0.9,1.7$, so the second observed sample is selected.
The mean is 0.333 and is not an observed sample. A medoid therefore is
\emph{not} a geometrical average shape. The denoised prototype is a separate
constructed shape combining compatible strokes; incompatible stroke counts
or order make naive vertex averaging invalid. The implementation's L1
consensus and cosine prior should not be silently substituted for one another.

\subsection{Feature extraction, connectedness and complete profile algorithm}
Slant comes from a shear search that concentrates ink into columns. Ascender
and descender extents, advances, lead spacing and word gaps are measured in
x-height units; robust summaries reduce isolated extraction errors. These
features are explicit geometric measurements, unlike learned CNN features.
Their meaning still depends on body-band and alignment quality.

Connectedness $c=N_{\rm cross}/\max(1,N_{\rm adj})$ counts eligible adjacent
non-space boundaries crossed by a traced centreline. Nine crossed boundaries
out of twelve give $c=0.75$. It is not fraction of ink pixels touching,
probability every new pair must join, or proof of original pen continuity.
A bad cut can manufacture a crossing, and detached crossbars can produce
false joins. If there are no eligible pairs the guarded formula returns zero;
that represents no observed crossings, not strong evidence of printing.

The full profile process is: pair source lines with labels; prepare the
reader input; force-align each label; estimate the body band/slant; reject
unusable alignments; refine one ordered boundary array; thin and trace the
full line; split chains at cuts; associate detached marks; normalise geometry;
record spacing and boundary crossings; cache extraction; reject weak lines;
filter each character's pool; compute prior and consensus; retain ranked
variants and representative forms; store robust style summaries and provenance.
Caching changes cost, not correctness, and stale cached preprocessing can
make a profile incompatible with current source. A usable glyph pool proves
only that the filters accepted it. It does not establish that all letters
or every writer's full style were captured.
